% cap_introducao.tex
\chapter{Introdução}
\label{chap:introducao}

A poluição atmosférica por material particulado fino (\textit{Particulate Matter}, PM2.5) permanece entre os problemas ambientais de maior impacto direto sobre a saúde pública. Por possuir diâmetro aerodinâmico inferior a 2,5 micrômetros, esse poluente consegue penetrar profundamente no sistema respiratório e está associado a desfechos respiratórios e cardiovasculares adversos \cite{who2021air,pope2006health,brook2010particulate}. A carga global de doença atribuível à poluição atmosférica por partículas finas reforça a importância de sistemas de monitoramento, prevenção e resposta \cite{cohen2017estimates}. No Brasil, episódios associados a queimadas e transporte atmosférico também ampliam a relevância do tema, com impactos mensuráveis sobre internações e mortalidade \cite{requia2021health,requia2022short}.

Nesse contexto, prever a concentração horária de PM2.5 com antecedência útil pode apoiar alertas, análise ambiental e planejamento operacional em regiões monitoradas. A previsão de curto prazo busca estimar valores futuros a partir de medições recentes e de variáveis auxiliares disponíveis no instante da previsão. Em uma formulação horária, o modelo deve reconhecer padrões de persistência, ciclos diários, mudanças graduais e eventos abruptos para produzir estimativas úteis para as próximas horas.

A tarefa, entretanto, é tecnicamente desafiadora. Séries ambientais reais combinam autocorrelação, sazonalidade, picos raros, mudanças de regime, ruído instrumental e lacunas de medição. Além disso, PM2.5 depende de condições atmosféricas, fontes de emissão, transporte de poluentes e relações com outros poluentes, nem sempre medidos de forma completa. Quando essas variáveis não estão disponíveis no momento da previsão, o modelo precisa operar com informação parcial, o que limita o desempenho possível mesmo com arquiteturas mais sofisticadas.

A literatura recente tem explorado modelos de aprendizado profundo para previsão de PM2.5, incluindo redes recorrentes, LSTM e arquiteturas com codificador, decodificador e atenção \cite{yang2023attention,zhao2021pm2,zhang2021deep}. Essas abordagens são atrativas porque conseguem modelar dependências temporais e relações não lineares. Ao mesmo tempo, revisões da área apontam dificuldade de comparação entre trabalhos, pois resultados dependem do conjunto de dados, horizonte de previsão, variáveis utilizadas e desenho de validação \cite{pm25_review_2024}. Em séries temporais, a separação cronológica e a avaliação fora da amostra também são parte central da validade experimental \cite{tashman2000out,bergmeir2012crossvalidation,hyndman2021forecasting}.

Este trabalho parte desse problema metodológico. Em vez de propor uma nova arquitetura como contribuição central, o TCC desenvolve um estudo experimental comparativo para previsão horária de PM2.5 na estação Sapo, pertencente ao recorte CMD (Conceição do Mato Dentro, MG) do conjunto de dados utilizado neste trabalho. A fonte bruta pertence ao monitoramento contínuo de qualidade do ar disponibilizado pela Fundação Estadual do Meio Ambiente de Minas Gerais \cite{feam2026dadosmonitoramento}. A escolha da estação é detalhada na metodologia, mas se apoia em três aspectos principais: cobertura do alvo PM2.5, disponibilidade pareada de PM10 e ocorrência suficiente de eventos de maior concentração.

\section{Motivação}

A motivação deste TCC é prática e metodológica. Do ponto de vista prático, previsões horárias de PM2.5 podem apoiar decisões de curto prazo em qualidade do ar, especialmente quando há risco de aumento de concentração nas próximas horas. Para isso, não basta obter uma métrica média aceitável: o modelo também precisa ser avaliado quanto à sua capacidade de preservar a amplitude de picos e de manter desempenho ao longo do horizonte de previsão.

Do ponto de vista metodológico, a comparação entre modelos só é informativa quando as decisões de dados e avaliação são controladas. Caso cada modelo use outro recorte temporal, outro conjunto de atributos, outra imputação ou outro teste, a comparação deixa de medir apenas a escolha de modelagem. Neste trabalho, um protocolo rastreável significa que a estação, o período, a tarefa, a preparação dos dados, os atributos disponíveis, a separação temporal e as métricas são explicitados e mantidos de forma comum sempre que possível. Assim, a análise busca responder não apenas qual modelo teve menor erro médio, mas também quais decisões explicam ganhos, falhas e compromissos observados.

A formulação principal usa uma janela de 48 horas passadas para prever as próximas 24 horas. Essa escolha é motivada por uma necessidade operacional de previsão do dia seguinte e pela presença de padrões diários em séries horárias de qualidade do ar. As 48 horas antecedentes permitem ao modelo observar dois ciclos diários recentes, enquanto o horizonte de 24 horas preserva uma tarefa de curto prazo suficientemente exigente para comparar degradação por horizonte. A justificativa empírica e os detalhes do protocolo são apresentados nos capítulos de fundamentação e metodologia.

\section{Problema de Pesquisa}

Para investigar o problema, o estudo compara modelos sequenciais e um modelo de referência tabular sob a mesma preparação de dados e o mesmo período cronológico de teste reservado. Esse modelo de referência é o XGBoost com saída múltipla, usado como comparação forte não sequencial para que os ganhos dos modelos recorrentes e Seq2Seq sejam interpretados com cautela.

A pergunta central deste trabalho pode ser formulada da seguinte forma:

\begin{quote}
Em uma tarefa horária de previsão de PM2.5 na estação Sapo, como diferentes escolhas de modelagem e preparação dos dados afetam o erro médio, os erros por horizonte e a capacidade de representar eventos de maior concentração?
\end{quote}

\section{Objetivos}

\subsection{Objetivo Geral}

Avaliar, sob um protocolo experimental rastreável e comparável, o desempenho de modelos de aprendizado de máquina e aprendizado profundo para previsão horária de múltiplos passos de PM2.5 na estação Sapo, distinguindo erro médio global, erros por horizonte e capacidade de representar regimes de maior concentração.

\subsection{Objetivos Específicos}

\begin{enumerate}
    \item Definir um protocolo rastreável de previsão horária de PM2.5 para a estação Sapo, com tarefa de múltiplos passos, entrada de 48 horas e horizonte de 24 horas.
    \item Consolidar uma base de dados com PM2.5, PM10, variáveis temporais, variáveis exógenas causalmente disponíveis, imputação controlada, normalização e exclusão de valores imputados do alvo na perda e nas métricas.
    \item Comparar LSTM, Seq2Seq e XGBoost multissaída, usando o XGBoost como modelo de referência tabular forte sob a mesma preparação de dados e o mesmo período cronológico de teste.
    \item Avaliar os resultados por métricas globais, métricas por horizonte e comportamento em picos/cauda, separando erro médio de capacidade de representar eventos de maior concentração.
    \item Documentar resultados negativos, limitações e compromissos de modelagem, especialmente no uso de atenção, perdas ponderadas, variáveis exógenas e estratégias voltadas a picos.
\end{enumerate}

\section{Organização do Trabalho}

Este documento está estruturado em seis capítulos. O Capítulo \ref{chap:introducao} apresenta o contexto, a motivação, o problema de pesquisa e os objetivos. O Capítulo \ref{chap:fundamentacao_teorica} reúne os conceitos necessários para fundamentar o trabalho, incluindo PM2.5, séries temporais, dados faltantes, modelos recorrentes, Seq2Seq, atenção, XGBoost e métricas de avaliação. O Capítulo \ref{chap:trabalhos_relacionados} posiciona o trabalho na literatura de previsão de qualidade do ar e discute lacunas de comparabilidade experimental. O Capítulo \ref{chap:materiais_metodos} descreve a metodologia do trabalho, incluindo dados, protocolo experimental, modelos, treinamento e métricas. O Capítulo \ref{chap:resultados} apresenta os resultados principais e as análises dos experimentos realizados. Por fim, o Capítulo \ref{chap:conclusao} sintetiza as conclusões, limitações e possibilidades de continuidade.
